\chapter*{Аннотация}
\addcontentsline{toc}{chapter}{Аннотация}

Работа посвящена применению методов машинного обучения и больших языковых моделей для финансового и операционного анализа малого производственного бизнеса. На примере предприятия, производящего кондитерские изделия и выпечку, построен воспроизводимый пайплайн обработки банковских и учётных данных: классификация транзакций (rule-based и LLM-подходы), управленческий cash flow, финансовый аудит денежных потоков и прогнозные модели спроса, возвратов и денежного потока. Показано, что при существенном обороте предприятие практически не формирует свободный операционный денежный поток, а ключевые риски связаны с концентрацией контрагентов и структурой возвратов по номенклатуре. Результаты интегрированы в собственную CRM-платформу, объединяющую загрузку данных, отчётность и прогнозные модели в одном интерфейсе, что превращает разовый аудит в регулярный управленческий инструмент.

\textbf{Ключевые слова:} машинное обучение, большие языковые модели, малый и средний бизнес, классификация транзакций, управленческий cash flow, прогнозирование спроса, CRM-платформа.

\chapter*{Abstract}
\addcontentsline{toc}{chapter}{Abstract}

This thesis examines the application of machine learning and large language models to financial and operational analysis in a small manufacturing business. Using a confectionery and bakery producer as a case study, a reproducible data pipeline is built for bank and accounting records: transaction classification (rule-based and LLM approaches), management cash flow statements, a financial audit of cash flows, and forecasting models for demand, returns, and cash flow. The analysis shows that despite substantial turnover, the company generates almost no free operating cash flow, with key risks stemming from customer concentration and the structure of product returns. The results are integrated into a custom-built CRM platform that combines data ingestion, reporting, and forecasting models in a single interface, turning a one-off audit into a recurring management tool.

\textbf{Keywords:} machine learning, large language models, small and medium-sized enterprises, transaction classification, management cash flow, demand forecasting, CRM platform.

\clearpage
